% CLAUDE REPORT::WCDERIV-0001::2026-11-14
\documentclass[11pt]{article}
\usepackage{amsmath,amssymb,amsthm}
\usepackage[margin=1in]{geometry}
\newtheorem{theorem}{Theorem}\newtheorem{corollary}{Corollary}\newtheorem{definition}{Definition}
\title{Worst-case minimum derivation cost $\overline W_R(N)$ and the weakest bridge\\
\small CLAUDE run WCDERIV-0001 (scheduled routine), 2026-11-14}
\date{}
\begin{document}\maketitle
\noindent\textbf{Status.} P versus NP remains \textbf{OPEN}. This run proves an exact
reformulation of the bridge (Theorem~1) and two corollaries that classify it. Nothing here
promotes any main-chain node. Living log:
\texttt{research/agents/claude/worst-case-derivation/LOG.md}.

\section{Objective (as assigned)}
\[
\overline W_R(N)=\max_{F\in\mathrm{UNSAT},\ |\langle F\rangle|\le N}\ \min_{\pi:F\vdash_R\bot}\operatorname{cost}_R(\pi),
\qquad
\textbf{(B$_R$)}\quad \mathrm{SAT}\in\mathrm P\ \Longrightarrow\ \overline W_R(N)\le \mathrm{poly}(N).
\]
Minimizing over all derivations with the optimal one supplied for free is exactly the
Cook--Reckhow proof-length measure, so it is a legitimate lower-bound strategy: any
algorithm that outputs a derivation pays at least its length. (KNOWN THEOREM, framing.)

\section{Corpus reviewed / duplication audit}
Surveyed the repository tree, \texttt{AGENTS.md}, the 2026-11-13 user-directed full-chain
handoff (target $W_{\rm abs}\le W_R\le c_M(n+T_M(n))^{r_M}$), recent TMR-1..5 handoffs,
the world ledger, and repository text matching ``worst-case derivation'' /
``minimum derivation cost'' (LSPV-1/348, 353, 365; LSPV-5/454--459). \textbf{Duplication:}
the corpus target $W_R(F_i)\le c_M(n_i+T_M(n_i))^{r_M}$ is the pointwise, per-algorithm form of
(B$_R$); this note analyses the weaker worst-case form and does not duplicate a proof
of either. \textbf{Contradiction:} none found. \textbf{Relative to root main:} no change to
root \texttt{main.tex} (this routine is not an LSPV owner); additive report only.

\section{Literature preflight (this run)}
\begin{itemize}
\item Cook--Reckhow (1979): a polynomially bounded proof system exists iff NP $=$ coNP.
\item Pich--Santhanam, \emph{Towards P$\ne$NP from Extended Frege lower bounds}
(arXiv:2312.08163), Prop.~1, and Arteche--Khaniki--Pich--Santhanam (ICALP 2024,
arXiv:2405.02232): proving unconditionally ``$S$ not p-bounded $\Rightarrow$ NP$\not\subseteq$P/poly''
for a concrete $S$ would already give NP$\not\subseteq$SIZE$[n^k]$ for every $k$. This is
the non-uniform analogue of (B$_R$) and is a \emph{formal barrier}.
\item Krajíček--Pudlák (1989): NE $=$ coNE implies an optimal proof system exists (existence only,
not for a fixed natural $R$). It is \emph{not} known that P $=$ NP makes EF polynomially bounded.
\item Haken 1985, Urquhart 1987, Ben-Sasson--Wigderson 2001: exponential resolution lower
bounds (PHP, Tseitin), so $\overline W_{\rm Res}(N)=2^{\Omega(N^{c})}$ unconditionally.
\end{itemize}

\section{Results}
\begin{definition}
For a deterministic polynomial-time machine $A$ and length $n$, let $\mathrm{Corr}_n(A)$ be
the propositional tautology (size $\mathrm{poly}(n)$, built from the computation tableau of $A$)
stating: for every $x\in\{0,1\}^n$ and every assignment $a$, if $A(x)$ rejects then $a$ does
not satisfy the formula encoded by $x$.
\end{definition}

\begin{theorem}[PROVED/EXACT; folklore-level]
Let $R$ be a Cook--Reckhow system that p-simulates Extended Frege, and measure $\overline W_R$
on 3-CNF refutations with $\operatorname{cost}=$ proof length. Then
\[
(\mathrm B_R)\iff\Bigl[\mathrm P=\mathrm{NP}\ \Rightarrow\ \exists\text{ poly-time }A\text{ deciding SAT with }
\mathrm{Corr}_n(A)\text{ having }\mathrm{poly}(n)\text{-size }R\text{-proofs}\Bigr].
\]
\end{theorem}
\begin{proof}
($\Leftarrow$) Given unsatisfiable $F$ with $|\langle F\rangle|=n$: substitute the constants $\langle F\rangle$ for $x$ in the
$R$-proof of $\mathrm{Corr}_n(A)$. EF (hence $R$) verifies the tableau of $A$ on these constants with
$\mathrm{poly}(n)$ extension steps, giving the rejection hypothesis. Modus ponens then gives
$\forall a\,\neg\mathrm{Sat}(\langle F\rangle,a)$. Turning that into a refutation of $F$ itself is a
polynomial-size EF step. So $\overline W_R(N)\le\mathrm{poly}(N)$.
($\Rightarrow$) Suppose P $=$ NP via some $A$, and suppose (B$_R$) holds, so $\overline W_R$ is polynomial. $\neg\mathrm{Corr}_n(A)$ is unsatisfiable
because $A$ is correct. Its Tseitin 3-CNF translation is unsatisfiable, has size $\mathrm{poly}(n)$, and has a polynomial
$R$-refutation. Since $R\ge$ EF, the extension variables let us convert that refutation into an $R$-proof of
$\mathrm{Corr}_n(A)$ of polynomial size.
\end{proof}

\noindent So, for every system at least as strong as EF, the bridge is exactly a
\emph{provable-correctness} statement. ``Run the fast algorithm and record the trace'' is not a
derivation until $R$ can certify the algorithm's correctness at each length. This is the precise
missing arrow. It matches the corpus warning that ``a correctness certificate for selected
algorithms is not the universal fixed-calculus transfer''.

\begin{corollary}[PROVED/EXACT]
If $\overline W_R$ is known to be superpolynomial, then (B$_R$) is equivalent to P $\ne$ NP.
This applies to resolution, tree-like systems, bounded-depth Frege, cutting planes and
polynomial calculus over any field, and to any multisign DAG-hypergraph clause-reduction
calculus that resolution p-simulates (merge, eliminate and subsume on signed hyperedges).
For these systems, proving (B$_R$) is a \emph{reformulation} (world-ledger category
``equivalent of P$\ne$NP''), not progress.
\end{corollary}
\begin{proof}
($\Rightarrow$) Contrapositive with the known lower bound. ($\Leftarrow$) If P $\ne$ NP, the
hypothesis of (B$_R$) is false, so (B$_R$) holds vacuously.
\end{proof}

\begin{corollary}[PROVED/EXACT]
For every $R$, (B$_R$) implies (P $\ne$ NP or NP $=$ coNP). Hence (B$_R$) plus any superpolynomial lower bound on $\overline W_R$ gives
P $\ne$ NP. This is the intended decomposition. For every $R$ studied so far, one of the two halves is open and
already at least as hard as a major open problem: for weak $R$ the bridge is equivalent to P $\ne$ NP (Corollary 1); for EF and stronger the lower bound is open.
The uniform analogue of the Pich--Santhanam barrier remains \textbf{OPEN}: we have not shown that proving (B$_R$) for a concrete
EF-strong $R$ forces a uniform breakthrough.
\end{corollary}

\section{Red-team computation (COMPUTATIONAL)}
Fixed reduction system $R_{\rm DP}$: Davis--Putnam elimination on the signed clause hypergraph,
with tautology and subsumption removal. The cost is the total number of clauses produced, minimized over all elimination orders by
subset dynamic programming; the order search is not charged. The value is an \emph{upper} bound on the true minimum, because subsumption can make
the reached state depend on the order. The control is GF(2) Gaussian elimination on the same odd-charge Tseitin
instances. Script \texttt{worst-case-derivation/dp\_vs\_gf2.py}; raw data
\texttt{2026-11-14\_dp\_vs\_gf2\_results.json}. Each graph runs as an independent chunk in under 10 s. Petersen
(15 variables) was skipped under the 30 s rule.
\begin{center}\begin{tabular}{lrrrr}
graph & vars & clauses & min DP cost (ub) & GF(2) row ops\\\hline
cycle$_{4,6,8,10}$ & 4--10 & 8--20 & 7, 11, 15, 19 & 3, 5, 7, 9\\
$K_4$ & 6 & 16 & 27 & 5\\
$K_{3,3}$ & 9 & 24 & 59 & 11\\
prism$_3$ & 9 & 24 & 47 & 9\\
cube$_3$ & 12 & 32 & 91 & 20\\
Wagner $V_8$ & 12 & 32 & 91 & 7
\end{tabular}\end{center}
On cycles (treewidth 2) the minimum DP cost is linear. It grows faster on higher-treewidth graphs,
consistent with the Ben-Sasson--Wigderson width/treewidth bounds, while GF(2) stays cubic. These sizes are tiny, so this table shows
no asymptotic behaviour. Its point is qualitative. Tseitin hardness of a clause-reduction calculus \textbf{fails to survive algebraic
reasoning}, so a hard family for (B$_R$) must be hard for the strongest $R$ under study, not
for the reduction rules one happens to prefer. (FAILED TRANSFER of Tseitin hardness to $R\ni$ GF(2) rules.)

\section{Barrier-transfer and PQ classification}
Barrier-transfer domain: logic and bounded arithmetic (reflection principles).
$R$ p-simulates a system $Q$ iff $R$ proves $Q$'s reflection principle with short proofs. That is the
\emph{directly transferable theorem} behind Theorem~1. PQ ledger item 6 (special
soundness): there is an ANALOGY/CANDIDATE that asks whether wrong rejections of $A$ can be \emph{extracted} as
explicit satisfying assignments. That would make $\mathrm{Corr}_n(A)$ a search-to-witness statement, but it
still needs a self-reducibility argument that is provable in $R$. Nothing transfers exactly.

\section{Adopted / rejected / retired; stale claims}
ADOPTED: free optimal sequence (minimum over derivations) as the cost measure. REJECTED as
progress: (B$_R$) for any $R$ with a known lower bound (Corollary~1). RETIRED: none. STALE CLAIMS:
none found in the files reviewed.

\section{Next targets per role}
\begin{enumerate}
\item Lead: try to prove or refute the \emph{uniform} Pich--Santhanam-type barrier for
(B$_{\rm EF}$).
\item Builder: fix $R=$ EF $+$ GF(2) $+$ cutting planes; write down the $\mathrm{Corr}_n$ tableau encoding with exact sizes.
\item Hypergraph/certificates: look for candidate families with no known short proofs in EF-strong
systems, such as random 3-CNF near threshold or the Razborov hardness candidates.
\item Auditor: charge the substitution and tableau steps of Theorem~1 exactly (degree $r_M$).
\item Red team: look for a natural $R$ in which (B$_R$) is \emph{false} under some
oracle or world (Impagliazzo world ledger), to rule out relativizing proofs of (B$_R$).
\end{enumerate}
P versus NP remains OPEN.
\end{document}
% END_OF_VERIFIED_REPORT::CLAUDE::WCDERIV-0001
